\chapter{Gemelo estático, discrepancia de fase e identificabilidad}
\label{ch:avances}

Este capítulo concentra la evidencia propia ya obtenida. Las cifras se reportan una
sola vez, junto con particiones, observables y límites de inferencia. Su propósito es
determinar qué se aprendió del canal estático y qué incertidumbre debe resolver la
experimentación física; no se presentan como resultados de sensado humano.

\section{Banco dc01 y condiciones de evaluación}

DICHASUS dc01 proporciona CFR MIMO-OFDM, posiciones y una escena compatible con trazado de rayos a \SI{3.438}{GHz} \citep{Euchner2023DICHASUSdcxx,DICHASUSdcxxWeb}. El banco contiene 10,000 posiciones; la partición original utiliza 5000 para entrenamiento, 100 para validación y 4900 para prueba. Los análisis posteriores añaden bloques espaciales con guarda para distinguir interpolación densa de generalización regional.

El canal medido y el simulado se comparan sobre soporte común. Potencia, RMS-DS, magnitud, coherencia y error angular responden preguntas distintas. Las correcciones que consultan el canal objetivo se etiquetan como oráculo o ajuste a posteriori; sólo un predictor entrenado fuera del test se considera transferible.

La unidad experimental es la posición o consulta. Antenas y subportadoras comparten geometría y no se tratan como réplicas independientes. Las cifras de este capítulo son evidencia dentro de dc01 y no validación entre edificios, dispositivos o sesiones.

\section{Fidelidad estática B0--B2}

B0 utiliza geometría y materiales nominales. B1 optimiza parámetros globales de materiales siguiendo calibración diferenciable \citep{Hoydis2024DiffRT}. B2 conserva el soporte de trazado y aprende parámetros efectivos dependientes de la posición de interacción; no es un campo RF implícito como NeRF$^2$ \citep{Zhao2023NeRF2}.

\begin{table}[htbp]
\centering\small
\caption{Resultados de prueba de la calibración estática B0--B2.}
\label{tab:b012_consolidated}
\begin{tabular}{lrrr}
\toprule
Métrica & B0 & B1 & B2\\
\midrule
Correlación de potencia $\uparrow$ & 0.699 & 0.736 & \textbf{0.871}\\
MAE de potencia [dB] $\downarrow$ & 6.29 & 3.82 & \textbf{2.22}\\
Correlación RMS-DS $\uparrow$ & 0.173 & 0.278 & \textbf{0.647}\\
MAE RMS-DS [ns] $\downarrow$ & 43.21 & 27.32 & \textbf{16.50}\\
NMSE mediana de magnitud [dB] $\downarrow$ & $-3.17$ & $-3.92$ & \textbf{$-4.46$}\\
Coherencia compleja bruta $\uparrow$ & 0.156 & 0.155 & 0.145\\
Error angular bruto [$^\circ$] $\downarrow$ & 89.93 & 89.95 & 89.96\\
\bottomrule
\end{tabular}
\end{table}

B2 mejora de forma clara potencia, magnitud y estructura de retardos, pero no la concordancia angular bruta. Los parámetros aprendidos se interpretan como efectivos: pueden compensar materiales, geometría, antenas o mecanismos ausentes. El primer hallazgo estructural es que una mejora macroscópica del canal no implica una mejora de fase compleja.

\section{Estructura de la discrepancia y componente transferible}

Los estudios históricos S4.2--S4.5 responden tres preguntas: qué estructura de baja dimensión existe, qué parte se transfiere sin observar el objetivo y qué mejora produce sobre el canal. Una aproximación operacional es
\begin{equation}
\Delta\phi_{i,u,a,k}\simeq
-2\pi\nu_k\Delta\tau_{{\rm eff},i,u}
+\psi_{i,u}+r^{\rm rel}_{i,u,a,k}\pmod{2\pi},
\label{eq:phase_contribution_factorization}
\end{equation}
donde $\Delta\tau_{\rm eff}$ es un retardo relativo, $\psi$ una dirección común al arreglo y $r^{\rm rel}$ estructura espacial/espectral restante. La expresión clasifica funciones inferenciales; no identifica causas físicas independientes.

El alineamiento oráculo muestra que una pendiente espectral común explica una fracción importante del error. Los retardos estimados correlacionan entre B0/B1/B2 (aproximadamente 0.886, 0.834 y 0.819). La transferencia entre grupos de antenas reduce el error y la interpolación entre frecuencias pares/impares es más favorable que extrapolar entre medias bandas. Un error de \SI{2.5}{ns} ya produce $22.5^\circ$ a \SI{25}{MHz} de separación.

La partición original posee distancia mediana de sólo \SI{4.83}{cm} al vecino de entrenamiento. En bloques reservados, el MAE medio de retardo es \SI{27.56}{ns} con posición, \SI{19.29}{ns} con descriptores RT y \SI{18.69}{ns} al combinar ambos. Los descriptores físicos conservan utilidad regional dentro de dc01, sin establecer todavía transferencia de escena.

\begin{table}[htbp]
\centering\small
\caption{Efecto del retardo predicho sobre el canal. El desfase constante por enlace se ajusta con la medición evaluada en todas las filas.}
\label{tab:phase_contribution_summary}
\begin{tabular}{llrr}
\toprule
Protocolo & Corrección & Error [$^\circ$] & Coherencia\\
\midrule
Denso & sólo CPO & 81.42 & 0.156\\
Denso & retardo predicho + CPO & \textbf{45.59} & \textbf{0.634}\\
Denso & retardo oráculo + CPO & 31.63 & 0.745\\
Por bloques & sólo CPO & 81.49 & 0.157\\
Por bloques & retardo predicho + CPO & \textbf{64.08} & \textbf{0.448}\\
Por bloques & retardo oráculo + CPO & 31.46 & 0.746\\
\bottomrule
\end{tabular}
\end{table}

La coherencia es invariante a una rotación constante por enlace; su aumento confirma el efecto de la corrección espectral. Sin embargo, el canal corregido todavía usa la medición evaluada para ajustar CPO. La contribución demostrada es el retardo transferible y su utilidad condicionada, no una predicción autónoma de fase absoluta.

\section{Fase común al arreglo}

S4.6 analiza 640,000 fases residuales en 10,000 capturas. La resultante global cercana a $0.0011$ y las resultantes marginales pequeñas no respaldan una fase absoluta fija y concentrada. Dentro de cada captura, la dirección común estimada reduce descriptivamente el error mediano de aproximadamente $89.0^\circ$ a $46.0^\circ$, pero se obtiene de las mismas antenas evaluadas.

La recuperación temporal produce correspondencia exacta de posiciones y permite comparar vecinos. Copiar $\psi$ desde el vecino espacial o temporal genera errores próximos a $90^\circ$ y no supera de forma clara los controles. El resultado excluye esas reglas locales bajo los intervalos examinados; no demuestra independencia ni determina una causa instrumental única.

La clasificación conceptual se expresa como
\begin{equation}
Z_\phi=\{Z_\phi^{\rm pred},Z_\phi^{\rm nuisance},Z_\phi^{\rm residual}\},
\label{eq:phase_latent_roles}
\end{equation}
donde el retardo efectivo ocupa el primer bloque de forma parcial, la fase común se asigna operacionalmente al segundo bajo el protocolo actual y la estructura relativa permanece en el tercero.

\section{S4.7-R: representaciones relativas e información espacial}
\label{sec:s47r_results}

S4.7-R evalúa si retirar grados de libertad de fase mejora la correspondencia espacial sin usar fase absoluta. El régimen es estático: no hay movimiento mecánico, referencia temporal independiente ni variable humana. La evaluación usa 500 consultas estratificadas y negativos locales, del mismo bloque y de otros bloques.

\begin{table}[htbp]
\centering\small
\caption{Discriminabilidad de $U(1)$ por arreglo antes y después del retardo predicho.}
\label{tab:s47_discriminability}
\begin{tabular}{lrrr}
\toprule
Métrica & RT bruto & RT con $\widehat\tau$ & Cambio\\
\midrule
Victoria global & 0.591333 & 0.795733 & 0.204400\\
Victoria local & 0.5200 & 0.5560 & 0.0360\\
Mismo bloque, no local & 0.5888 & 0.8872 & 0.2984\\
Entre bloques & 0.6652 & 0.9440 & 0.2788\\
Exactitud \emph{top-1} & 0.048 & 0.156 & 0.108\\
\bottomrule
\end{tabular}
\end{table}

Los intervalos pareados son
\[
\mathrm{IC95\%}_{\rm global}=[0.1801,0.2281],\qquad
\mathrm{IC95\%}_{\rm local}=[-0.0016,0.0732].
\]
La mejora local no es concluyente. El producto espectral con separación 64 ofrece información regional complementaria. La resolución conjunta sostenida aparece entre $1$ y $2\lambda$ (aproximadamente \SIrange{0.087}{0.174}{m}); es un umbral empírico de discriminabilidad, no un límite físico universal.

La fusión convexa se ajusta sin usar test y asigna peso $0.551$ a $U(1)$ y $0.449$ al producto espectral; magnitud y el control R4 reciben peso nulo o despreciable.

\begin{table}[htbp]
\centering\small
\caption{Evaluación congelada de la fusión de S4.7-R.}
\label{tab:s47_frozen_fusion}
\begin{tabular}{lrrr}
\toprule
Métrica & $U(1)$ & Fusión & Cambio\\
\midrule
Victoria global & 0.795733 & 0.807333 & 0.011600\\
Victoria local & 0.5560 & 0.5560 & 0.0000\\
Mismo bloque, no local & 0.8872 & 0.8928 & 0.0056\\
Entre bloques & 0.9440 & 0.9732 & 0.0292\\
Exactitud \emph{top-1} & 0.156 & 0.148 & $-0.008$\\
\bottomrule
\end{tabular}
\end{table}

La mejora global tiene IC95\% $[0.0031,0.0207]$; entre bloques, el intervalo es
$[0.0192,0.0408]$. No se demuestra mejora local, \emph{top-1} o MRR. La representación relativa mejora discriminabilidad espacial estática, pero no demuestra fidelidad diferencial ni sensing.

\section{S4.7-Ua: predictibilidad circular de la fase común}

S4.7-Ua compara $p(\psi\mid x_{\rm RT})$ con baselines circulares incondicionales mediante NLL y particiones fuera de muestra. La dirección condicionada no muestra una mejora robusta; retirar el término $\beta$ no cambia la conclusión; el efecto local restringido a $\kappa>0$ permanece alrededor de cero; y una escala oráculo RT no produce una ganancia material. El mínimo observado, $\Delta\mathrm{NLL}\approx-2\times10^{-6}$, es despreciable.

\begin{quote}
No se encontró evidencia de predictibilidad circular transferible de la fase común a partir de los descriptores evaluados.
\end{quote}

En consecuencia, $\psi\in Z_\phi^{\rm nuisance}$ bajo dc01/B0, las covariables y particiones evaluadas. El cierre es operacional. No autoriza afirmar que $\psi$ sea universalmente aleatoria o que ninguna referencia futura pueda volverla predecible.

\section{S4.7-Ub.0: resolubilidad del modelo por trayectorias}

La reconstrucción de la suma de caminos alcanza error aproximado de $2.93\times10^{-7}$, lo que verifica la implementación de síntesis. La inversión de coeficientes revela, sin embargo, una diferencia severa:
\begin{table}[htbp]
\centering\small
\caption{Diagnóstico inicial de la matriz pathwise.}
\label{tab:pathwise_smoke}
\begin{tabular}{lrrr}
\toprule
Caso & Rango & $\kappa(G)$ & Coherencia máxima\\
\midrule
Tx0 & 189/189 & $2.79\times10^2$ & 0.997073\\
Tx1 & 178/178 & $1.45\times10^6$ & 0.999974\\
\bottomrule
\end{tabular}
\end{table}

Las ecuaciones normales deterioran notablemente Tx1 porque cuadran la condición. Un ridge exploratorio reduce varianza a costa de un sesgo demasiado alto y no constituye una solución. Con $B\approx\SI{50}{MHz}$, la escala $1/B\approx\SI{20}{ns}$ ayuda a explicar firmas de caminos casi indistinguibles. El hallazgo metodológico es preciso: rango algebraico completo no implica identificabilidad robusta de coeficientes pathwise.

El siguiente análisis estudiará valores singulares, rango efectivo, coherencias, separación de retardos y agrupación de modos no resolubles. No se presenta todavía un posterior independiente por trayectoria como resultado científico.

\section{Evidencia dinámica complementaria de d010}

DICHASUS d010 aporta una trayectoria temporal del transmisor observada por infraestructura distribuida \citep{DICHASUSd010Web}. Su variación confirma que la sensibilidad depende de geometría y enlaces, y resulta útil para auditar continuidad y particiones temporales. No constituye sensing humano con terminales fijos:
\begin{equation}
\frac{\partial H}{\partial p_{\rm Tx}}
\ne
\frac{\partial H}{\partial q_{\rm humano}}.
\label{eq:d010_derivative_distinction}
\end{equation}
En el segundo caso coexisten caminos estáticos y dinámicos, cambios de dispersión y oclusiones corporales. d010 se conserva como evidencia externa complementaria, no como sustituto de E1--E3.

\section{Síntesis científica y consecuencias}
\label{sec:synthesis}

La calibración estática caracteriza una parte importante de la discrepancia, y el retardo efectivo exhibe estructura transferible limitada. La fase común no muestra predictibilidad condicionada útil en S4.7-Ua. Las representaciones relativas preservan información espacial, pero todavía se desconoce cuánto Jacobiano físico conservan. La inversión pathwise está restringida por resolubilidad y condición numérica.

Estos resultados estrechan la hipótesis doctoral: el siguiente avance no depende de añadir capacidad neuronal a dc01, sino de medir piso instrumental y perturbaciones controladas. E0 y E1 deben determinar si las representaciones seleccionadas reducen deriva sin destruir $\Delta H$ y si el gemelo reproduce $J_q$. Sólo entonces la evidencia estática podrá conectarse rigurosamente con fisiología y actividad humana.

\section{Trazabilidad histórica}
\label{sec:trazabilidad_avances}

La nomenclatura S4.2--S4.7 se conserva para enlazar registros de experimentación. B0--B2 corresponde a calibración estática; S4.2--S4.5 a estructura y transferencia del retardo; S4.6 a fase común; S4.7-R a representaciones relativas; S4.7-Ua a predictibilidad circular; S4.7-Ub.0 a resolubilidad pathwise. Los nombres de programas, hashes, parámetros completos y estructura de artefactos se mantienen en la documentación reproducible del repositorio, fuera de la narrativa científica principal.
